\documentclass[11pt]{article}
\usepackage[margin=1in]{geometry}
\usepackage{amsmath,amssymb,amsthm}
\usepackage[T1]{fontenc}
\usepackage{lmodern}

\title{Literature-implied answer for an equal-weight Parseval fusion-frame subcase}
\author{Run fa\_banach\_001, lane 17}
\date{2026-07-03}

\begin{document}
\maketitle

\paragraph{Status.}
\texttt{literature\_implied\_answer (partial subcase)}.

\paragraph{Source question.}
The source paper is arXiv:2303.01202, L. Kohldorfer, P. Balazs,
P. Casazza, S. Heineken, C. Hollomey, P. Morillas, and M. Shamsabadi,
\emph{A Survey of Fusion Frames in Hilbert Spaces}. In the introductory
``There Be Dragons'' list and again in Remark \texttt{rank1fusion}, the paper
states that there are currently no necessary and sufficient conditions for the
existence of Parseval fusion frames for arbitrary subspace dimensions. The
example given there is that two two-dimensional subspaces of $\mathbb R^3$
cannot form a Parseval fusion frame for any choice of weights.

\paragraph{Supporting theorem.}
The supporting paper is arXiv:1112.3060, M. Bownik, K. Luoto, and
E. Richmond, \emph{A combinatorial characterization of tight fusion frames}.
Its Problem 1 asks for a characterization of sequences
$(L_1,\ldots,L_K)$ for which there exist orthogonal projections
$P_1,\ldots,P_K$ on $\mathbb R^N$ with
\[
  P_1+\cdots+P_K=\alpha I_N,
  \qquad \operatorname{rank}P_i=L_i.
\]
Theorem \texttt{th:combchar} gives the necessary and sufficient criterion:
if $M=\sum_i L_i$ and $\alpha=M/N$, then such projections exist if and only if
\[
  c\bigl((N^{L_1}),\ldots,(N^{L_K});(M^N)\bigr)\ne 0,
\]
where this is the corresponding Littlewood--Richardson coefficient. The same
paper notes that this removes the equal-rank limitation of the earlier
Casazza-Fickus-Mixon-Wang-Zhou result.

\paragraph{Identification.}
A finite Parseval fusion frame with common weight $w>0$ and subspace ranks
$L_i$ satisfies
\[
  w^2(P_1+\cdots+P_K)=I_N.
\]
Thus it is exactly the same data as an unweighted tight fusion frame with
frame bound $\alpha=w^{-2}$. Conversely, any unweighted tight fusion frame
$P_1+\cdots+P_K=\alpha I_N$ becomes Parseval after assigning the common weight
$w=\alpha^{-1/2}$ to all subspaces. Therefore Bownik-Luoto-Richmond give a
necessary and sufficient condition for the finite equal-weight Parseval
dimension problem.

\paragraph{Scope limitation.}
This packet does not claim a full answer to the arbitrary-weight problem
$\sum_i w_i^2P_i=I_N$ with unconstrained, nonconstant weights, nor to the
fixed-subspace weight-scalability problem studied in later work. It records
only the equal-weight finite-dimensional subcase as already covered by known
literature after the above identification.

\paragraph{Local evidence.}
No local PDFs were present for arXiv:2303.01202 or arXiv:1112.3060. The
evidence used here is the parsed TeX source at
\texttt{data/parsed/arxiv\_sources/2303.01202/Main.tex} and
\texttt{data/parsed/arxiv\_sources/1112.3060/honeycombfinal.tex}.

\end{document}
